\section{Hermite coefficients of $\Xi$}\label{tr:sec:hermite}

We expand $\Xi$ in Hermite functions at a small scale, following \cite{P26hermite}, whose definitions are written out here. The kernel is the P\'olya kernel $\Phi$ of \eqref{eq:ch2:Phi}, even and positive with $\Xi(t)=\int_{\mathbb R}\Phi(u)e^{itu}du$; with it the recomputed numbers of the note agree, except where stated, whereas the de Bruijn form $\Phi_{\rm dBN}$ of Section~\ref{ch:dbn} puts the zeros at $2\gamma$. The kernel satisfies $\int|\Phi(u)|(1+|u|)^Ne^{Bu^2+C|u|}du<\infty$ for all $N,B,C\ge0$ (it decays like $\exp(-\pi e^{2|u|})$), and this is all that is used.

\subsection{Definitions}

Let $H_n$ be the Hermite polynomial of Section~\ref{ht:sec:jensen} (physicists' normalisation, the one that gives $H_{2d}$ in the Hermite identity there), and
\[
h_n(x)=\pi^{-1/4}(2^nn!)^{-1/2}H_n(x)e^{-x^2/2},
\]
the orthonormal Hermite functions. They satisfy $h_n''+(2n+1-x^2)h_n=0$, $\int h_n(x)e^{i\xi x}dx=\sqrt{2\pi}\,i^nh_n(\xi)$, and $h_{2m}(0)=(-1)^m\pi^{-1/4}\sqrt{(2m)!}/(2^mm!)$. For a scale $a>0$ put
\begin{gather*}
c_n(a)=\sqrt a\int_{\mathbb R}\Phi(u)h_n(au)\,du,\qquad t_m=t_m(a)=a\sqrt{4m+1},\\
b_m(a)=\sqrt{2\pi}\,(-1)^mc_{2m}(a)=\int_{\mathbb R}\Xi(t)\,\frac{h_{2m}(t/a)}{\sqrt a}\,dt.
\end{gather*}
The second form of $b_m$ follows from the Fourier relation with $t=ax$. Since $\Phi$ is even, $c_n(a)=0$ for odd $n$, and $b_m(a)$ is the coefficient of the orthonormal function $h_{2m}(t/a)/\sqrt a$ in the expansion of $\Xi$. The number $t_m$ is the turning point of $h_{2m}(t/a)$: the equation $h_{2m}''+(4m+1-x^2)h_{2m}=0$ is oscillatory for $|x|<\sqrt{4m+1}$, that is for $|t|<t_m$.

With $H_t$ the de Bruijn Newman function of Section~\ref{ch:dbn}, $8H_t(2\tau)=\int_{\mathbb R}\Phi(v)e^{tv^2/4}\cos(\tau v)\,dv$, so the Gaussian smoothing
\[
\begin{aligned}G_a(\tau)&=\int_{\mathbb R}\Phi(u)e^{-a^2u^2/2}\cos(\tau u)\,du=8H_{-2a^2}(2\tau)\\&=\frac1{a\sqrt{2\pi}}\int_{\mathbb R}\Xi(s)e^{-(\tau-s)^2/(2a^2)}ds\end{aligned}
\]
is $\Xi$ smoothed by a Gaussian of variance $a^2$ ($e^{-a^2u^2/2}$ is the characteristic function of $N(0,a^2)$), the backward heat flow of the kernel for the time $t=-2a^2$.

\begin{theorem}[Hermite coefficients and the heat flow]\label{tr:thm:heat}
For every $a>0$,
\[
\sum_{m\ge0}b_m(a)\frac{z^{2m}}{\sqrt{(2m)!}}=\sqrt{2\pi}\,\pi^{-1/4}\sqrt a\;e^{z^2/2}\,G_a\big(\sqrt2\,a z\big),
\]
and with $\nu_k(a)=\int\Phi(u)e^{-a^2u^2/2}u^{2k}du>0$,
\[
b_m(a)=\sqrt{2\pi}\,\pi^{-1/4}\sqrt a\,\sqrt{(2m)!}\sum_{k=0}^m\frac{(-1)^k2^ka^{2k}\nu_k(a)}{(2k)!\,2^{m-k}(m-k)!}.
\]
\end{theorem}

\begin{proof}
The generating function
\[
\sum_nh_n(x)w^n/\sqrt{n!}=\pi^{-1/4}e^{-x^2/2+\sqrt2xw-w^2/2}
\]
integrated against $\sqrt a\,\Phi(u)$ at $x=au$ gives
\[
\sum_nc_n(a)w^n/\sqrt{n!}=\pi^{-1/4}\sqrt a\,e^{-w^2/2}\int\Phi(u)e^{-a^2u^2/2}e^{\sqrt2auw}du.
\]
On bounded sets of $w$ the integrand and its derivatives are bounded by a constant times $|\Phi(u)|e^{-a^2u^2/2+C|u|}$, which is integrable, so the coefficients may be extracted termwise. Put $w=\allowbreak iz$: the odd terms vanish by evenness, $(iz)^{2m}=\allowbreak (-1)^mz^{2m}$ supplies the sign in $b_m$, and $e^{i\sqrt2azu}$ integrates to $G_a(\sqrt2az)$ by evenness. The moment formula follows by multiplying the series
\[
G_a(\sqrt2az)=\sum_k(-1)^k2^ka^{2k}\nu_kz^{2k}/(2k)!
\]
by
\[
e^{z^2/2}=\sum_jz^{2j}/(2^jj!).
\]
Each $\nu_k>0$ because $\Phi>0$.
\end{proof}

The moment formula in $60$ digit arithmetic agrees with the quadrature value of $b_m$ for $m\le15$ at $a=3$ and $a=1.5$ with relative differences $6.7\times10^{-55}$ and $3.1\times10^{-59}$, and the two quadratures (steps $1/64$ and $1/128$) differ by $3.6\times10^{-57}$ and $1.5\times10^{-60}$ (computed). Also $b_3(1.5)=0.320600426659651$ from $\Phi$ and from the integral of $\Xi$ (computed).

\begin{remark}\label{tr:rem:rt}
By Rodgers and Tao \cite{RodgersTao}, $\Lambda\ge0$, so $H_t$ has nonreal zeros for every $t<0$, and therefore $G_a$ is not real rooted for any $a>0$, since $t=-2a^2<0$ (proved). Since $e^{z^2/2}$ never vanishes, $e^{z^2/2}G_a(\sqrt2az)$ has exactly the zeros of $G_a(\sqrt2az)$. The note's statement that the coefficients of this product nevertheless recover the real zeros of $\Xi$ is not about zeros of the product; what is true is that the signs of $b_m(a)$ follow those of $\Xi$ at the points $t_m$ away from the zeros as $a\to0$ (Theorem~\ref{tr:thm:herm}).
\end{remark}

\section{The turning point asymptotic}\label{tr:sec:saddle}

The note proposes, as Proposition 1, that $b_m(a)=\sqrt{2a}\,m^{-1/4}\,\Xi(a\sqrt{4m+1})(1+O(a^2))$ as $a\to0$ with $t=a\sqrt{4m+1}$ fixed, and derives it by a saddle point argument on the circle $|z|=\sqrt{2m}$ in Theorem~\ref{tr:thm:heat}. The leading constant is right: the saddles are at $\theta=0,\pi$ with the exponent $m-2m\theta^2+O(m\theta^3)$, the Gaussian in $\theta$ has variance $1/(4m)$, the imaginary shift of the argument of $G_a$ has variance exactly $a^2$, and $\mathbb EG_a(t+iaY)=\Xi(t)$ for $Y\sim N(0,1)$ by Fubini, so the average of $G_a$ over the shifts undoes the smoothing; Stirling's formula and the Gaussian integral give $\sqrt{2a}\,m^{-1/4}$. But the error $O(a^2)$ is not derived (only the shift of the sampling point from $2a\sqrt m$ to $a\sqrt{4m+1}$ is), the real displacement $-a\sqrt m\,\theta^2$ and the cubic phase are not estimated, and the contribution off the saddles is not bounded. The statement as written is false at the zeros of $\Xi$ (Proposition~\ref{tr:prop:zeros}). The theorem below, proved from the Hermite differential equation, replaces it with an additive estimate, uniform in $m$, with the exact amplitude.

Put
\begin{gather*}
A_m(a)=\sqrt{2\pi a}\,(-1)^mh_{2m}(0)=\sqrt2\,\pi^{1/4}\sqrt a\,\frac{\sqrt{(2m)!}}{2^mm!}>0 ,\\
M_k(a_0)=\int_{\mathbb R}|\Phi(u)||u|^ke^{a_0^2u^2/2}du<\infty .
\end{gather*}

\begin{theorem}[Additive asymptotic]\label{tr:thm:herm}
For every $a_0>0$, $0<a\le a_0$ and integer $m\ge0$,
\begin{gather*}
\Big|\frac{b_m(a)}{A_m(a)}-\Xi(t_m)\Big|\le\frac{M_4(a_0)}{12}\,a^4,\\
\frac{b_m(a)}{A_m(a)}=\Xi(t_m)+a^4Q(t_m)+R_m(a),\qquad|R_m(a)|\le\frac{M_8(a_0)}{672}\,a^8,
\end{gather*}
with $t_m=a\sqrt{4m+1}$ and the even entire function
\begin{gather*}
Q(t)=\frac{\Xi'''(t)}{6t}-\frac{\Xi''(t)}{4t^2}+\frac{\Xi'(t)}{4t^3}\quad(t\ne0),\\
Q(0)=\frac{\Xi^{(4)}(0)}{12}=\frac1{12}\int_{\mathbb R}\Phi(u)u^4du .
\end{gather*}
\end{theorem}

\begin{proof}
Put $y_m(u)=h_{2m}(au)/h_{2m}(0)$. Then $b_m/A_m=\int\Phi(u)y_m(u)du$, and the Hermite equation gives $y_m''+(t_m^2-a^4u^2)y_m=0$, $y_m(0)=1$, $y_m'(0)=0$. For real $t$ and complex $z$ let $y_{z,t}$ be the even solution of $y''+(t^2-zu^2)y=0$ with these initial values. For $u\ge0$, $y_{z,t}=\cos tu+z\,L_ty_{z,t}$ with $(L_tf)(u)=\int_0^u\frac{\sin t(u-v)}tv^2f(v)dv$ (the kernel is $u-v$ at $t=0$). Since $|\sin(t(u-v))/t|\le u-v$, the Volterra series is majorised by the solution of $w=1+|z|Kw$ with $(Kf)(u)=\int_0^u(u-v)v^2f(v)dv$. The function $v=e^{\sqrt{|z|}u^2/2}$ satisfies $v''=(\sqrt{|z|}+|z|u^2)v\ge|z|u^2v$, so $v\ge1+|z|Kv$, hence $w\le v$ and
\[
|y_{z,t}(u)|\le e^{\sqrt{|z|}u^2/2}\quad(u\in\mathbb R,\ t\in\mathbb R).
\]
Iterating the integral equation once and twice, and using $K1=u^4/12$ and $K^21=u^8/672$, gives $|y_{z,t}-\cos tu|\le\frac{|z|u^4}{12}e^{\sqrt{|z|}u^2/2}$ and $|y_{z,t}-\cos tu-zL_t\cos(t\,\cdot)|\le\frac{|z|^2u^8}{672}e^{\sqrt{|z|}u^2/2}$, with constants independent of $t$. Take $z=a^4$, $t=t_m$, $a\le a_0$ and integrate against $\Phi$: $\int\Phi\cos(t_mu)du=\Xi(t_m)$, and the first inequality follows. The function $f_t=(\frac{u^3}{6t}-\frac u{4t^3})\sin tu+\frac{u^2}{4t^2}\cos tu$ solves $f''+t^2f=u^2\cos tu$ with $f(0)=f'(0)=0$, so $f_t=L_t\cos(t\,\cdot)$, and differentiating the Fourier integral of $\Xi$ gives $\int\Phi f_t=Q(t)$. At $t=0$, $f_0=u^4/12$.
\end{proof}

For the amplitude, Stirling's expansion gives $A_m(a)=\sqrt{2a}\,m^{-1/4}\big(1-\frac1{16m}+\frac1{512m^2}+O(m^{-3})\big)$. Hence for $0<\varepsilon\le t=t_m\le T$ (so that $m\to\infty$ uniformly as $a\to0$)
\begin{equation}\label{tr:eq:amp}
\frac{b_m(a)}{\sqrt{2a}\,m^{-1/4}}=\Big(1-\frac1{16m}+O(m^{-2})\Big)\big(\Xi(t)+O(a^4)\big),
\end{equation}
an additive statement, uniform even if the interval contains zeros. At a fixed $t>0$ with $\Xi(t)\ne0$ one may divide, and then $b_m(a)=\sqrt{2a}\,m^{-1/4}\,\Xi(t)(1+O_t(a^2))$, which is the formula of the note; the leading correction is $-1/(16m)=-a^2/(4t^2)+O(a^4)$ and comes from the amplitude and not from the saddle. A uniform relative error holds on a compact set $K\subset(0,\infty)$ with $\inf_K|\Xi|>0$. With the sampling point as the variable, $n=m+\frac14=t^2/(4a^2)$, the amplitude is $A_m=\frac{2a}{\sqrt t}\big(1-\frac{a^4}{8t^4}+O(a^8)\big)$ and
\[
b_m(a)=\frac{2a}{\sqrt t}\Big\{\Xi(t)+a^4\Big[Q(t)-\frac{\Xi(t)}{8t^4}\Big]+O_{\varepsilon,T}(a^8)\Big\},
\]
with no term of order $a^2$ or $a^6$; this form is not uniform near $t=0$, where Theorem~\ref{tr:thm:herm} with the exact $A_m$ is the right one. At $m=0$ one has $t_0=a$ and $b_0(a)=A_0(a)\int\Phi(u)e^{-a^2u^2/2}du>0$, and the formula with $m^{-1/4}$ is undefined; at fixed $m\ge1$, $b_m/A_m\to\Xi(0)>0$ as $a\to0$. Even for $m\ge1$ the relative formula with the prefactor $\sqrt{2a}$ does not extend to $t=0$: for $m=1$, $b_1(a)/(\sqrt{2a}\,\Xi(a\sqrt5))\to\pi^{1/4}/\sqrt2=0.9414$ (proved, from $A_1=\pi^{1/4}\sqrt a$).

\begin{proposition}\label{tr:prop:zeros}
At every positive real zero $\gamma$ of $\Xi$, of any multiplicity, $b_m(a_m)\ne0$ for all large $m$, where $a_m=\gamma/\sqrt{4m+1}$. If $Q(\gamma)\ne0$ then $b_m(a_m)=\frac{2Q(\gamma)}{\sqrt\gamma}a_m^5+O(a_m^9)$. The relative formula of the note therefore fails at the zeros of $\Xi$, where it would give $b_m(a_m)=0$.
\end{proposition}

\begin{proof}
For fixed real $t$ put $F_t(z)=\int\Phi(u)y_{z,t}(u)du$. By the bound of the proof of Theorem~\ref{tr:thm:herm} it is entire in $z$, and $F_t(0)=\Xi(t)$. For $t>0$ and $z=t^4$ the solution is $y=e^{-t^2u^2/2}$, so $F_t(t^4)=\int\Phi e^{-t^2u^2/2}du>0$. At $t=\gamma$, $F_\gamma(0)=0$ but $F_\gamma\not\equiv0$, so $F_\gamma(z)\ne0$ for small $z\ne0$; and $b_m(a_m)=A_m(a_m)F_\gamma(a_m^4)$. If $Q(\gamma)\ne0$ the Taylor expansion at $z=0$ gives the leading term, with $A_m(a_m)\sim2a_m/\sqrt\gamma$.
\end{proof}

\section{The first negative coefficient}\label{tr:sec:first}

Let $m^*(a)=\min\{m\ge0:b_m(a)<0\}$, and let $\gamma_1=14.134725\ldots$ be the first positive zero of $\Xi$. It is simple (the first $1.5\times10^9$ zeros are simple and on the line \cite{LRW86}), so $\Xi>0$ on $[0,\gamma_1)$ and $\Xi<0$ just to the right of $\gamma_1$.

\begin{corollary}\label{tr:cor:first}
For every $0<\delta<\gamma_1$ and $a$ small, $b_m(a)>0$ for all $m$ with $t_m\le\gamma_1-\delta$, the index $m^*(a)$ exists, and $\lim_{a\to0}4a^2m^*(a)=\gamma_1^2$. If $\gamma_a$ is the unique zero near $\gamma_1$ of $t\mapsto\int\Phi(u)y_{a^4,t}(u)du$, then $\gamma_a=\gamma_1-a^4Q(\gamma_1)/\Xi'(\gamma_1)+O(a^8)$ and $m^*(a)=\big\lfloor(\gamma_a^2/a^2-1)/4\big\rfloor+1$ for small $a$.
\end{corollary}

\begin{proof}
Let $c_\delta=\min_{[0,\gamma_1-\delta]}\Xi>0$. By Theorem~\ref{tr:thm:herm}, $b_m/A_m\ge c_\delta-Ca^4>0$ when $t_m\le\gamma_1-\delta$ and $a$ is small, so $m^*(a)$, where it exists, has $t_{m^*}>\gamma_1-\delta$ and $4a^2m^*=t_{m^*}^2-a^2>(\gamma_1-\delta)^2-a^2$, which gives $\liminf4a^2m^*\ge\gamma_1^2$. For the upper bound pick $s>\gamma_1$ with $\Xi(s)<0$ and $m(a)=\lfloor(s^2/a^2-1)/4\rfloor$; then $t_{m(a)}\to s$, the additive error tends to zero, so $b_{m(a)}<0$ and $\limsup4a^2m^*\le s^2$; let $s\downarrow\gamma_1$. For the exact formula, $F(z,t)=\int\Phi y_{z,t}$ is jointly analytic near $(0,\gamma_1)$ (the kernel bound $|\sin(tw)/t|\le we^{|\Im t|w}$ gives an integrable majorant), $\partial_zF(0,t)=Q(t)$ and $\partial_tF(0,t)=\Xi'(t)\ne0$, so the analytic implicit function theorem gives a zero $\rho(z)=\gamma_1-zQ(\gamma_1)/\Xi'(\gamma_1)+O(z^2)$ and $\gamma_a=\rho(a^4)$. For small $a$ the function $F(a^4,\cdot)$ is positive up to $\gamma_a$ and negative just after it, and the grid has spacing $O(a^2)$ there, so the first negative index is the first grid point strictly after $\gamma_a$.
\end{proof}

The note's hypothesis, an error uniform on $[0,\gamma_1-\delta]$, gives only the lower bound: the limit needs negativity of $b_m$ just above $\gamma_1$, so a uniform estimate on an interval reaching beyond $\gamma_1$.

The note states the sharper law $m^*(a)=\lceil(\gamma_1^2/a^2-1)/4\rceil$ for $a\le1.5$. Table~\ref{tr:tab:first} (computed, second run at step $h/2$ the same) agrees with it at four scales, but it cannot hold for every $a$. The threshold inside the rounding is $(\gamma_a^2/a^2-1)/4$, which differs from $(\gamma_1^2/a^2-1)/4$ by $-\gamma_1Q(\gamma_1)a^2/(2\Xi'(\gamma_1))+O(a^6)$, and $\gamma_a\ne\gamma_1$ for small $a$ (otherwise $F_{\gamma_1}$ would vanish identically, against $F_{\gamma_1}(\gamma_1^4)>0$), so $\gamma_a-\gamma_1=c\,a^{4p}+o(a^{4p})$ with $c\ne0$ and $p\ge1$. If $c>0$, at $a=\gamma_1/\sqrt{4m+1}$ the ceiling of the note gives $m$ and the first negative coefficient is $m+1$ for large $m$; if $c<0$ a slightly smaller $a$ gives the opposite. Such scales exist arbitrarily close to $0$, because the displacement of the root is much smaller than the spacing $O(a^2)$ of the grid. At $a=2$ the formula gives $12.24$ and the first negative index is $12$, so the ceiling formula, which would give $13$, fails there, although $a=2$ is not in the asymptotic regime; and the reading of $\gamma_1$ to about $10^{-2}$ from $m^*(1/2)$ is only to $0.035$ (the value $m^*(\frac12)=200$ places $\gamma_1$ in $(14.116,14.151]$).

\begin{table}[h]\centering\small
\begin{tabular}{lccccc}
\toprule
$a$ & $0.5$ & $0.75$ & $1$ & $1.5$ & $2$\\
\midrule
$(\gamma_1^2/a^2-1)/4$ & $199.54$ & $88.55$ & $49.70$ & $21.95$ & $12.24$\\
first negative $m$ & $200$ & $89$ & $50$ & $22$ & $12$\\
\bottomrule
\end{tabular}
\caption{The first negative Hermite coefficient of $\Xi$ (computed).}\label{tr:tab:first}
\end{table}

\section{Counting zeros from $\Phi$ alone}\label{tr:sec:count}

\subsection{The ratio table}

Table~\ref{tr:tab:ratio} gives the ratio $b_m(a)\big/\big(\sqrt{2a}\,m^{-1/4}\,\Xi(a\sqrt{4m+1})\big)$ at three scales (computed; the note does not say which $m$, and the indices $m=9,36,81,144,182$ at $a=0.5$, $4,16,36,64,81$ at $a=0.75$ and $2,9,20,36,45$ at $a=1$ reproduce every printed digit; the sampling points are $t_m=3.04,6.02,9.01,12.01,13.50$, $3.09,6.05,9.03,12.02,13.52$ and $3.00,6.08,9.00,12.04,13.45$). The last column is the error at $t\approx9$ divided by $a^2$ (the note printed $0.0032$ for the first). The table is explained by the amplitude: the printed values are close to $A_m/(\sqrt{2a}\,m^{-1/4})=1-\frac1{16m}+\cdots$ (for example $1-1/144=0.9931$ at $a=0.5$, $m=9$), and the term $a^4Q(t)/\Xi(t)$ of Theorem~\ref{tr:thm:herm} is the rest. The coefficients range over four orders of magnitude. At larger $a$ and greater height the ratio is no longer near $1$; the note says it is between $1.5$ and $9$ at $a=4$ for $100<t<200$ and that no $H_\lambda$ with $\lambda\in[-8,3]$ fits better, which could not be reproduced (the ratio is not defined near the zeros of $\Xi$, and it varies in sign over the window).

\begin{table}[h]\centering\small
\begin{tabular}{ccccccc}
\toprule
$a$ & $t\approx3$ & $t\approx6$ & $t\approx9$ & $t\approx12$ & $t\approx13.5$ & error at $t\approx9$ over $a^2$\\
\midrule
$0.50$ & $0.9931$ & $0.9983$ & $0.9992$ & $0.9993$ & $0.9986$ & $0.0033$\\
$0.75$ & $0.9847$ & $0.9961$ & $0.9980$ & $0.9976$ & $0.9937$ & $0.0036$\\
$1.00$ & $0.9699$ & $0.9931$ & $0.9959$ & $0.9938$ & $0.9830$ & $0.0041$\\
\bottomrule
\end{tabular}
\caption{The ratio of $b_m(a)$ to $\sqrt{2a}\,m^{-1/4}\,\Xi(a\sqrt{4m+1})$ (computed).}\label{tr:tab:ratio}
\end{table}

\subsection{Method and the count}

The integrand $\Phi(u)h_{2m}(au)$ is not entire: it is analytic in $|\Im u|<\pi/4$ and the lines $\Im u=\pm\pi/4$ are a natural boundary (\eqref{eq:ch2:Phi} and the text after it), but it decays like $\exp(-\pi e^{2|u|})$ and the trapezoid rule converges geometrically. For a function analytic in a strip $|\Im u|\le d$ with $M_d=\int(|f(x+id)|+|f(x-id)|)dx<\infty$ the error of the infinite trapezoid rule of step $h$ is at most $M_d/(e^{2\pi d/h}-1)$ (proved, by shifting the Fourier integral to the edges of the strip and summing two geometric series), with $d<\pi/4$; the constants depend on $a$ and $m$, so the note's statement that the transform of the integrand is negligible beyond $2T$ and that $h=2\pi/(3T+100)$ suffices is not proved. The count was recomputed with this step (the note has about $655$ nodes at $T=500$), the Hermite functions by the three term recurrence for all orders at once, and a working precision of the order of $0.34\,T+40$ digits (the coefficients near height $T$ have size $e^{-\pi T/4}$), and the positions of the sign changes were compared with a run at half the step and more digits (computed). The note's column of digits does not follow its own rule and is omitted.

\begin{table}[h]\centering\small
\fitwidth{\begin{tabular}{rrrl}
\toprule
$T$ & $N(T)$ & coefficients & sign changes of $b_m$ by scale $a$\\
\midrule
$100$ & $29$ & $625$ at $a=2$ & $a=2,3,4,5$: $29$, $29$, $29$, $29$\\
$200$ & $79$ & $625$ at $a=4$ & $a=4,5$: $79$, $79$; $a=6,7$: $78$, $78$; $a=8$: $77$\\
$500$ & $269$ & $1736$ at $a=6$ & $a=6$: $269$; $7$: $268$; $8$: $267$; $9$: $265$; $10$: $263$; $14$: $245$; $18$: $192$\\
$1000$ & $649$ & $6945$ at $a=6$ & $a=5,6$: $648$, $648$ (edge lag)\\
$1004$ & $652$ & $7000$ at $a=6$ & $a=6$: $651$\\
\bottomrule
\end{tabular}}
\caption{Number of sign changes of $b_m(a)$ for $t_m(a)\le T$ against the zero count $N(T)$ (computed). The coefficient counts of the note ($627$, $1738$, $6947$, $7002$) were two too large.}\label{tr:tab:count}
\end{table}

\begin{remark}[Observation]\label{tr:obs:count}
For $a$ up to about $5$ to $6$ the number of sign changes of $(b_m(a))_{t_m\le T}$ equals $N(T)$ at $T=100,200,500$, and at $T=1000$ it equals $N(T)$ except for the last zero within about one unit of the top of the computed range. In every run above the threshold the count was smaller than $N(T)$.
\end{remark}

This is an observation (computed, floating point; the positions agree in the second run at $T=100$, $T=200$ and $T=500$, $a=6$, the other rows being single runs), not a theorem: neither a universal one sided inequality nor an association between sign changes of $b_m$ and zeros of zeta is proved. The zero counts are those of zeta ($N(T)=29,79,269,649,652$, mpmath). The equality is a statement about counts and not about signs: at $T=200$, $a=4$ the sign of $b_m$ agrees with the sign of $\Xi(t_m)$ at $81.8$ percent of the indices ($c=1$; $81.9$ percent at $c=0$), and sampling at $a\sqrt{4m+c}$ for $c\in[0,3]$ moves this between $79.7$ and $81.9$ percent (computed).

\subsection{The edge lag at $T=1000$}

At $T=1000$ both $a=5$ and $a=6$ give $648$ sign changes against $N(1000)=649$. The last zero is $\gamma_{649}=999.79$, $0.3$ below the top of the range. Extending the computation at $a=6$ to $t\le1004$ the next sign changes occur between consecutive $t_m$ at $(1000.85,1000.92)$, $(1002.07,1002.14)$ and $(1003.36,1003.44)$, for the zeros $999.79$, $1001.35$ and $1002.40$ (computed), and $\gamma_{652}=1003.27$ again falls inside the lag at the new edge; the note gives the sign changes at $1000.89$, $1002.11$, $1003.40$, which are inside these intervals. The sign changes trail the zeros by about one unit at this height, and the count must be read on $[0,T-\delta]$ with a lag $\delta(a,T)$. The note adds that at low height the turning points sit slightly below the zeros and at $T=1000$ slightly above, so that $\delta$ changes sign with height; this is an observation that Theorem~\ref{tr:thm:herm} does not yet explain.

\begin{proposition}[What a verified count proves]\label{tr:prop:verified}
Let $\widetilde b_m/A_m$ be computed values of the normalised coefficients with rigorous errors $e_m$, and let $E_m$ be a rigorous bound for $|b_m/A_m-\Xi(t_m)|$ (for example $M_4(a_0)a^4/12$ from Theorem~\ref{tr:thm:herm}). If $|\widetilde b_m/A_m|>e_m+E_m$ then $\operatorname{sgn}\widetilde b_m=\operatorname{sgn}\Xi(t_m)$. Suppose there are $S$ sign changes at pairs of adjacent sampling points below $T$, all with this property. Then $\Xi$ has $S$ zeros in disjoint intervals, and if an independent rigorous method gives $N(T)\le S$, where $N(T)$ is the number of zeros $t$ of $\Xi$ with $0<\Re t<T$, counted with multiplicity (Turing's method, Section~\ref{sec:ch5:turing}), all those zeros are real and simple.
\end{proposition}

\begin{proof}
The triangle inequality transfers the signs, the intermediate value theorem gives one zero of $\Xi$ in each of the $S$ disjoint intervals between adjacent sampling points, and the upper bound, counted with multiplicity, leaves no room for another zero or a multiple zero.
\end{proof}

A sign change of $b_m$ certifies nothing by itself, since $b_m$ is an integral of $\Xi$, not a value. The chain of the note, a rigorous version of the asymptotic, a lower bound for $|\Xi|$ at the sampling points and Turing's method, is Proposition~\ref{tr:prop:verified} with Theorem~\ref{tr:thm:herm} supplying the first step with explicit constants; the sampling values that are zero or uncertain must be enclosed separately. High precision and agreement after halving the step are diagnostics, not these inequalities, so the counts of Table~\ref{tr:tab:count} are computed, not verified.

\section{The best scale and the run lengths}\label{tr:sec:scale}

Two quantities compete. The grid $t_m(a)$ has spacing $dt_m/dm=2a^2/t_m$ (exactly $4a^2/(t_{m+1}+t_m)$ between neighbours), so a small $a$ resolves close zeros; the number of coefficients is $M=T^2/(4a^2)+1$, so a large $a$ is cheap. The note says that the binding constraint is not resolution, because at the largest successful scale the grid spacing is only $0.12$ to $0.17$ of the mean gap $2\pi/\log(T/2\pi)$. The recomputed ratios of $2a^2/T$ to the mean gap are $0.22$ ($T=100$, $a=5$), $0.14$ ($T=200$, $a=5$), $0.10$ ($T=500$, $a=6$) and $0.06$ ($T=1000$, $a=6$) (computed), so the note's range is not reproduced; and the inference does not follow in any case, since two zeros arbitrarily close together can fit between two sampling points whatever the mean gap (the function $(t-c)^2-\varepsilon^2$ is positive at both ends of a grid interval for small $\varepsilon$). The tables do not separate a missed grid crossing from a distortion of the asymptotic; what is observed is that the counts fall short above $a\approx5$ to $6$ and that the usable threshold rises slowly with $T$, from at least $5$ at $T=100$ to between $6$ and $7$ at $T=500$. With $a$ near the threshold there are about $6.5$ coefficients per zero at $T=500$ and $10.7$ at $T=1000$ ($1736/269$ and $6945/649$), each costing $O(T)$ nodes at $0.34\,T$ digits. As a verification method this is far slower than Odlyzko and Sch\"onhage \cite{OS88}; its value is structural, since every quantity is an explicit integral of the positive kernel $\Phi$.

The note also reads the sign pattern through Good and Gaskins \cite{GoodGaskins80}: with $\Xi(t)=-\frac12(t^2+\frac14)\pi^{-1/4}|\Gamma(\frac14+\frac{it}2)|Z(t)$ for the Hardy function $Z$ and the Riemann Siegel phase $\vartheta'(t)=\frac12\log(t/2\pi)+O(t^{-2})$, the sign of $b_m$ would rotate by $\theta_m=\vartheta'(t_m)\,dt_m/dm=a^2\log(t_m/2\pi)/t_m$ per step in $m$, an effective complex ratio, with mean sign run $\pi/\theta_m=\pi t/(a^2\log(t/2\pi))$. The rotation law is not derived: the real function $Z$ is not determined by its phase derivative, no recurrence or complex ratio for $b_m$ has been derived, and Theorem~\ref{tr:thm:herm} is an asymptotic at bounded height as $a\to0$ and does not apply at fixed $a=4$ and increasing height (the attribution to Good and Gaskins is the note's and was not checked).
% UNCHECKED: Good and Gaskins 1980, attribution of the effective complex ratio is the source's and was not checked.
What can be said is an averaged statement.

\begin{proposition}[Average run length]\label{tr:prop:run}
Let $N(T)$ count the zeros $t$ of $\Xi$ with $0<\Re t<T$, with multiplicity. Assume the zeros $t$ of $\Xi$ with $T\le\Re t\le T+H$ are real and simple and that the sign changes of $(b_m(a))$ with $t_m\in[T,T+H]$ correspond one to one to them. Then the mean number of coefficient steps per sign change is the number $TH/(2a^2)+O(H^2/a^2+1)$ of grid points in the window divided by $N(T+H)-N(T)=\frac{H}{2\pi}\log\frac T{2\pi}+O(H^2/T+\log T)$, and for fixed $a$, $H\to\infty$ and $H=o(T)$ the ratio is asymptotic to $\pi T/(a^2\log(T/2\pi))$.
\end{proposition}

Table~\ref{tr:tab:run} tests this conditional average, with the predicted values at the window midpoints $60,100,140,180$ and $a=4$ (arithmetic). The note's observed values $5.29$, $7.19$, $8.89$, $10.50$ were not reproduced, because "mean run" is not defined in the note (edge blocks, which window); with the number of indices in the window divided by the number of sign blocks at $T=200$, $a=4$, the observed values of the table (computed) are consistent with the average and do not support the rotation law.

\begin{table}[h]\centering\small
\fitwidth{\begin{tabular}{lcccc}
\toprule
window in $t$ ($a=4$) & $40$ to $80$ & $80$ to $120$ & $120$ to $160$ & $160$ to $200$\\
\midrule
predicted $\pi t/(a^2\log(t/2\pi))$ & $5.22$ & $7.10$ & $8.86$ & $10.53$\\
observed mean run (indices over sign blocks) & $5.00$ & $6.94$ & $8.33$ & $10.23$\\
\bottomrule
\end{tabular}}
\caption{Mean sign runs of $b_m(4)$ (computed).}\label{tr:tab:run}
\end{table}

\section{The Thorin shadow}\label{tr:sec:shadow}

The note concludes that, by the turning point asymptotic, Hermite coefficients of a real even function sample the function at the points $a\sqrt{4m+1}$, so that sign definiteness of the coefficients is asymptotically equivalent to positivity of the function on the real line, and that a Thorin or GGC positivity statement for the reciprocal kernel has a Hermite form that is no stronger and no weaker. This does not follow and is not true as stated. Let $f(t)=(\sin t/t)^2$, which is even, entire and nonnegative on the real axis, with $f(t)=\frac14\int_{-2}^2(2-|u|)e^{itu}du$ and a nonnegative compactly supported kernel. The proof of Theorem~\ref{tr:thm:herm} applies to it. At $t=\pi$ one has $f'(\pi)=0$, $f''(\pi)=2/\pi^2$, $f'''(\pi)=-12/\pi^3$, so $Q_f(\pi)=-5/(2\pi^4)<0$, and along $a=\pi/\sqrt{4m+1}$ the normalised coefficient is $a^4Q_f(\pi)+O(a^8)<0$ for all large $m$: a nonnegative function can have negative small scale Hermite coefficients at every large $m$ along that sequence (proved). The correct statements are two. For a kernel of the class above, if $\Xi$ (or any function of the class) is strictly positive on a compact interval, then every coefficient with $t_m$ in the interval is positive once $a$ is small (the positive minimum and the additive error, with $A_m>0$). Conversely, if all these coefficients are nonnegative for every small $a$, the function is nonnegative on the interval (choose grid points converging to any $t$ and pass to the limit in the normalised formula; at $t=0$ use $m=0$). Neither characterises the Thorin or GGC property of a reciprocal kernel. The Hermite coefficients add computability; no result saying more is known here.

\section{Notes and open problems}\label{tr:sec:trnotes}

Nothing computed here is certified in ball arithmetic. The Lehmer pair near $t=7005$ that the note mentions consists of the zeros $6709$ and $6710$ at $7005.062866$ and $7005.100565$, a gap of $0.0377$ (computed).



\begin{problem}\label{tr:prob:gram}
Prove that Gram anomalies with the margin $B_m(\gamma)$ of Proposition~\ref{tr:prop:gram} occur infinitely often for $L(s,\chi_4)$, adapting the arguments known for $\zeta$ \cite{Trudgian11b}. Together with Proposition~\ref{tr:prop:gram} this would make Corollary~\ref{tr:cor:nogo} unconditional at each order $m$ with $\eta_m>0$.
\end{problem}

\begin{problem}\label{tr:prob:F8}
Prove that $F_8$, the symmetrised hard cutoff with eight atoms, has only critical zeros, by combining the phase of $g_8$ at low height with the modulus inequality and a Newton polygon bound at large height. The analogues for zeta are Theorem~\ref{rz:thm:Yn} and Corollary~\ref{rz:cor:Ssmall}.
\end{problem}

\begin{problem}\label{tr:prob:largest}
Determine the largest $n$ for which the symmetrised mean corrected truncation has only critical zeros in the strip, and decide whether the first failure is related to the first Gram anomaly at $138.75$.
\end{problem}

\begin{problem}\label{tr:prob:remainder}
Find a remainder $R_n$, not a cumulant match, whose drift $\eta\delta_m$ has a sign independent of the Gram parity. By Theorem~\ref{tr:thm:phase} this needs a correction whose phase cancels $\vartheta_\beta(\gamma)$, which suggests building $R_n$ from the functional equation and not from the tail.
\end{problem}

\begin{problem}\label{tr:prob:zeta}
Carry out the analysis of orders $m\ge1$ for zeta with the lattice of Section~\ref{ch:rz4}, where $m=0$ is Theorem~\ref{rz:thm:kappaexact}, and compare the Gram parity of the corrected phase with the classical failures of Gram's law for zeta (the first near the Gram point $g_{126}\approx282.455$, where $Z(g_{126})=-0.0276$).
\end{problem}

\begin{problem}\label{tr:prob:explicit}
Evaluate the constants of Theorem~\ref{tr:thm:herm} explicitly (the moments $M_4(a_0)$ and $M_8(a_0)$ and a bound for $Q$) and use Proposition~\ref{tr:prop:verified} with Turing's method to count the zeros of $\Xi$ up to a given height from the coefficients alone, including the edge lag $\delta(a,T)$.
\end{problem}

\begin{problem}\label{tr:prob:amp}
Explain the distortion of the counts at large $a$, and find a modified coefficient sequence, perhaps with $e^{z^2/2}$ replaced by another multiplier in Theorem~\ref{tr:thm:heat}, that keeps the count exact at larger $a$ and so cuts the number $M=T^2/(4a^2)$ of coefficients.
\end{problem}

\begin{problem}\label{tr:prob:lehmer}
Push the count to $T=2000$ and to the Lehmer pair near $t=7005$, where a sign change count is the natural place for a close pair to hide.
\end{problem}

\begin{problem}\label{tr:prob:shadow}
Find a Hermite statement that is equivalent to a Thorin or GGC positivity property of the reciprocal kernel, or show that none exists.
\end{problem}
